\documentclass[../../../../main.tex]{subfiles}
\begin{document}

\begin{flushright}
\footnotesize\textit{【自動文字起こし・要確認】}
\end{flushright}

[解] $\sin\theta = S, \cos\theta = C$ と書く。$\left(0 < \theta < \frac{\pi}{2}\right)$
\begin{enumerate}
\item[(1)] $0 < \theta < \frac{\pi}{2}$ から $0 < 2\theta < \pi, \quad 0 < 3\theta < \frac{3}{2}\pi$ より
\begin{align*}
\sin 3\theta &= \sin 2\theta \Leftrightarrow 3\theta = 2\theta,\ 3\theta = \pi - 2\theta \\
&\Leftrightarrow \theta = \frac{\pi}{5} \quad \left( \because 0 < \theta < \frac{\pi}{2} \right)
\end{align*}

\item[(2)] $0 < \theta < \frac{\pi}{2}$ に注意して
\begin{equation*}
(\text{与式}) \Leftrightarrow 4C^2 - 2mC - n - 1 = 0 \quad \cdots
\end{equation*}
よりこれが解を持つ条件は、左辺 $f(C)$ $(0 < C < 1)$、判別式 $D$ として
$f(0) = -(n+1) < 0$ に注意すると、$f(1) > 0 \Leftrightarrow 3-2m-n > 0$
である。$m, n$ が非負整数であることから、
\begin{equation*}
(m, n) = \underset{\text{\textcircled{1}}}{(0, 0)}, \underset{\text{\textcircled{2}}}{(0, 1)}, \underset{\text{\textcircled{3}}}{(0, 2)}, \underset{\text{\textcircled{4}}}{(1, 0)}
\end{equation*}
である。
\begin{itemize}
\item[$1^\circ$の時]
$f(C) = 4C^2 - 1 = 0 \quad \therefore C = \frac{1}{2} \quad \left(\theta = \frac{\pi}{3}\right)$
\item[$2^\circ$の時]
$f(C) = 4C^2 - 2 = 0 \quad \therefore C = \frac{\sqrt{2}}{2} \quad \therefore \theta = \frac{\pi}{4}$
\item[$3^\circ$の時]
$f(C) = 4C^2 - 3 = 0 \quad \therefore C = \frac{\sqrt{3}}{2}, \theta = \frac{\pi}{6}$
\item[$4^\circ$の時]
$(\text{与式}) \Leftrightarrow \sin 3\theta = \sin 2\theta$ だから (1) から $\theta = \frac{\pi}{5}$
\end{itemize}
以上に $0 < C < 1$, $0 < \theta < \frac{\pi}{2}$ を用いた。従って解は
\begin{equation*}
(m, n, \theta) = \left(0, 0, \frac{\pi}{3}\right), \left(0, 1, \frac{\pi}{4}\right), \left(0, 2, \frac{\pi}{6}\right), \left(1, 0, \frac{\pi}{5}\right)
\end{equation*}
\end{enumerate}

\end{document}
